% %%%%%%%%%%%%%%%%%%%%%%%%%%%%%%%%%%%%%%%%%
% % LaTeX Template
% % Version 2.0 - Elaborato ASDI
% %
% % This template originates from:
% % http://www.LaTeXTemplates.com
% %
% % Authors:
% % Rocco Lo Russo (https://github.com/roccolr)
% %%%%%%%%%%%%%%%%%%%%%%%%%%%%%%%%%%%%%%%%%
%
% Indice dei comandi e degli ambienti definiti qui (esempi d'uso in capitoli/):
%
%   INTESTAZIONE     \corso, \docente, \annoaccademico, \studente{Nome}{Matricola}
%   BOZZA            \todo{...}            (nascosto con \bozzafalse in main.tex)
%   CODICE           \begin{vhdl}[opzioni] ... \end{vhdl}
%                    \begin{xdc}[opzioni]  ... \end{xdc}
%                    \begin{tcl}[opzioni]  ... \end{tcl}
%                    \vhdlfile[opzioni]{percorso}{didascalia}{etichetta}
%                    \xdcfile[opzioni]{percorso}{didascalia}{etichetta}
%                    \sig{nome}   \ent{nome}   (segnali ed entity nel testo)
%   RIQUADRI         traccia, scelta, osservazione, attenzione, risultato
%                    (+ commandline, file, question, warn, info della v1.0)
%   INTERFACCE       \begin{porte} \porta{nome}{dir}{larghezza}{descrizione} \end{porte}
%   FIGURE           \screenshot[larghezza]{file}{didascalia}{etichetta}
%   SCHEMI TIKZ      stili: blocco, mux, reg, fsmstate, seg, bus, ctrl
%   SEZIONI STANDARD \progetto \implementazione \simulazione \sintesiboard \timinganalysis
%
%----------------------------------------------------------------------------------------


%----------------------------------------------------------------------------------------
%	PACKAGES AND OTHER DOCUMENT CONFIGURATIONS
%----------------------------------------------------------------------------------------

\usepackage[utf8]{inputenc} % Required for inputting international characters
\usepackage[T1]{fontenc} % Output font encoding for international characters
\usepackage[italian]{babel}

\usepackage{amsmath,amsfonts,stmaryrd,amssymb} % Math packages

\usepackage[ddmmyyyy]{datetime}
\usepackage{enumerate} % Custom item numbers for enumerations
\usepackage{enumitem}

\usepackage[ruled,linesnumbered,italiano]{algorithm2e} % Algorithms (parole chiave in italiano)

\usepackage{xcolor}
\usepackage{graphicx}
\graphicspath{{figure/}{figure/screenshot/}{../pics/}}

\usepackage{booktabs}  % Tabelle professionali (\toprule, \midrule, \bottomrule)
\usepackage{tabularx}
\usepackage{array}
\usepackage{multirow}
\usepackage{float}
\usepackage{caption}
\usepackage{subcaption}
\captionsetup{font=small, labelfont=bf, format=hang}

\usepackage[per-mode=symbol]{siunitx} % Unità di misura: \SI{10}{\nano\second}, \SI{100}{\mega\hertz}
\sisetup{output-decimal-marker={,}}

\usepackage[framemethod=tikz]{mdframed} % Allows defining custom boxed/framed environments

\usepackage{tikz}
\usetikzlibrary{positioning, arrows.meta, automata, shapes.geometric, shapes.misc, calc, fit, backgrounds}

\usepackage{listings} % File listings, with syntax highlighting

\usepackage{fancyhdr}
\usepackage{titlesec}

\usepackage[hidelinks]{hyperref}

%----------------------------------------------------------------------------------------
%	DOCUMENT MARGINS
%----------------------------------------------------------------------------------------

\usepackage{geometry} % Required for adjusting page dimensions and margins

\geometry{
	paper=a4paper, % Paper size, change to letterpaper for US letter size
	top=2.5cm, % Top margin
	bottom=3cm, % Bottom margin
	left=2.5cm, % Left margin
	right=2.5cm, % Right margin
	headheight=14pt, % Header height
	footskip=1.5cm, % Space from the bottom margin to the baseline of the footer
	headsep=1.2cm, % Space from the top margin to the baseline of the header
	%showframe, % Uncomment to show how the type block is set on the page
}

%----------------------------------------------------------------------------------------
%	FONTS
%----------------------------------------------------------------------------------------

\usepackage{XCharter} % Use the XCharter fonts
\usepackage[varqu,varl,scaled=0.95]{inconsolata} % Monospazio con grassetto per il codice
\usepackage[scaled=0.9]{helvet}  % Sans-serif per titoli e riquadri (XCharter non ne ha uno)

%----------------------------------------------------------------------------------------
%	COLORI
%----------------------------------------------------------------------------------------

\definecolor{codebg}{gray}{0.97}
\definecolor{accento}{HTML}{1F4E79}     % blu scuro: titoli, scelte progettuali
\definecolor{osservcol}{HTML}{2E7D32}   % verde: osservazioni
\definecolor{attcol}{HTML}{C62828}      % rosso: attenzione
\definecolor{riscol}{HTML}{6A1B9A}      % viola: risultati
\definecolor{traccol}{gray}{0.45}       % grigio: testo della traccia

%----------------------------------------------------------------------------------------
%	DATI DEL DOCUMENTO
%----------------------------------------------------------------------------------------

\newcommand{\corso}[1]{\def\thecorso{#1}}
\newcommand{\docente}[1]{\def\thedocente{#1}}
\newcommand{\annoaccademico}[1]{\def\theannoaccademico{#1}}
\def\thestudenti{}
% \studente{Nome Cognome}{Matricola}: può essere ripetuto per ogni componente del gruppo
\newcommand{\studente}[2]{\expandafter\def\expandafter\thestudenti\expandafter{\thestudenti #1 & \texttt{#2} \\}}

%----------------------------------------------------------------------------------------
%	MODALITÀ BOZZA
%----------------------------------------------------------------------------------------

% \todo{testo} evidenzia in rosso le parti ancora da scrivere.
% Con \bozzafalse in main.tex i \todo scompaiono dalla versione finale.
\newif\ifbozza
\bozzatrue
\newcommand{\todo}[1]{\ifbozza{\color{attcol}\sffamily\small[\textbf{TODO:} #1]}\fi}

%----------------------------------------------------------------------------------------
%	TITOLI, INTESTAZIONI
%----------------------------------------------------------------------------------------

\titleformat{\chapter}[display]
	{\normalfont\sffamily\color{accento}}
	{\large\MakeUppercase{\chaptertitlename}~\thechapter}
	{0.3em}
	{\Huge\bfseries}
	[\vspace{0.4em}{\color{accento}\titlerule[1pt]}]
\titlespacing*{\chapter}{0pt}{-1.5cm}{1.2cm}

\titleformat{\section}{\normalfont\Large\bfseries\sffamily\color{accento}}{\thesection}{0.8em}{}
\titleformat{\subsection}{\normalfont\large\bfseries\sffamily}{\thesubsection}{0.8em}{}
\titleformat{\subsubsection}{\normalfont\normalsize\bfseries\sffamily\color{black!75}}{}{0em}{}

\setcounter{secnumdepth}{2}  % numerati fino a \subsection (Esercizio 1.1)
\setcounter{tocdepth}{2}

% Sezioni standard richieste dal template dell'elaborato (non numerate, non nell'indice)
\newcommand{\progetto}{\subsubsection*{Progetto e architettura}}
\newcommand{\implementazione}{\subsubsection*{Implementazione}}
\newcommand{\simulazione}{\subsubsection*{Simulazione}}
\newcommand{\sintesiboard}{\subsubsection*{Sintesi su board di sviluppo}}
\newcommand{\timinganalysis}{\subsubsection*{Timing analysis}}

\pagestyle{fancy}
\fancyhf{}
\renewcommand{\chaptermark}[1]{\markboth{\thechapter.\ #1}{}}
\renewcommand{\sectionmark}[1]{\markright{\thesection\ #1}}
\fancyhead[L]{\small\sffamily\color{black!60}\leftmark}
\fancyhead[R]{\small\sffamily\color{black!60}\rightmark}
\fancyfoot[C]{\small\sffamily\thepage}
\renewcommand{\headrulewidth}{0.4pt}
\fancypagestyle{plain}{\fancyhf{}\fancyfoot[C]{\small\sffamily\thepage}\renewcommand{\headrulewidth}{0pt}}

%----------------------------------------------------------------------------------------
%	LISTINGS: STILI DI BASE
%----------------------------------------------------------------------------------------

\renewcommand{\lstlistingname}{Listato}
\renewcommand{\lstlistlistingname}{Elenco dei listati}

% Caratteri accentati nei listati (listings non gestisce UTF-8 da solo)
\lstset{
  literate=
    {à}{{\`a}}1 {è}{{\`e}}1 {é}{{\'e}}1 {ì}{{\`i}}1 {ò}{{\`o}}1 {ù}{{\`u}}1
    {À}{{\`A}}1 {È}{{\`E}}1 {É}{{\'E}}1 {Ì}{{\`I}}1 {Ò}{{\`O}}1 {Ù}{{\`U}}1
    {→}{{$\rightarrow$}}1 {←}{{$\leftarrow$}}1,
}

\lstdefinestyle{base}{
  basicstyle=\ttfamily\small\color{black},
  numbers=left,
  numberstyle=\scriptsize\color{gray},
  stepnumber=1,
  numbersep=6pt,
  frame=single,
  rulecolor=\color{black!30},
  backgroundcolor=\color{codebg},
  framesep=5pt,
  xleftmargin=20pt,
  framexleftmargin=15pt,
  breaklines=true,
  postbreak=\mbox{\textcolor{gray}{$\hookrightarrow$}\space},
  showstringspaces=false,
  tabsize=4,
  columns=flexible,
  keepspaces=true,
  captionpos=t,
  abovecaptionskip=4pt,
  keywordstyle=\color{purple}\bfseries,
  commentstyle=\color{teal},
  stringstyle=\color{orange},
  identifierstyle=\color{black},
  emphstyle=\color{blue}\bfseries,
}

% Stile C (quello originale della v1.0)
\lstdefinestyle{c}{
  style=base,
  language=C,
  emph={size_t, uint32_t, uint64_t, TAS, RTE},
  morekeywords={uint8_t, uint16_t, uint32_t, uint64_t, int8_t, int16_t, int32_t, int64_t, bool, true, false},
  morecomment=[l]{//},
  morecomment=[s]{/*}{*/},
}
\lstset{style=c}

\lstdefinelanguage{RISCAsm}{
  morekeywords={
    ADDQ, SUM, ANDI, MOVE,MOVEA, MOVEM, DC, DS, sra,
    lw, sw, li, la, mv,
    BEQ, BNE, JMP, CMP, jalr, ret,
    NOP, lui, RTE, RTS,
    ecall, ebreak, ORG, EQU, JSR, CLR, TAS
  },
  sensitive=true,
  morecomment=[l]{\*},
  morestring=[b]",
}

%----------------------------------------------------------------------------------------
%	LISTINGS: VHDL, XDC, TCL
%----------------------------------------------------------------------------------------

\lstdefinestyle{vhdl}{
  style=base,
  language=VHDL,
  sensitive=false,
  morekeywords={
    rising_edge, falling_edge, others, generate, generic, map, port,
    downto, to, when, else, elsif, select, with, case, is, of, range,
    wait, until, for, loop, report, severity, assert, note, warning, error, failure
  },
  % tipi e funzioni di libreria in blu
  emph={
    std_logic, std_logic_vector, std_ulogic, std_ulogic_vector,
    unsigned, signed, integer, natural, positive, boolean, bit, bit_vector,
    time, string, character,
    to_unsigned, to_signed, to_integer, resize, shift_left, shift_right,
    rotate_left, rotate_right, std_logic_1164, numeric_std, ieee, work
  },
  emphstyle=\color{blue}\bfseries,
}

\lstdefinestyle{xdc}{
  style=base,
  language=tcl,
  morekeywords={set_property, get_ports, get_cells, get_nets, create_clock, current_design,
    create_generated_clock, set_input_delay, set_output_delay, set_false_path},
  emph={PACKAGE_PIN, IOSTANDARD, LVCMOS33, LVCMOS18, CONFIG_VOLTAGE, CFGBVS, VCCO},
  emphstyle=\color{blue}\bfseries,
  morecomment=[l]{\#},
}

\lstdefinestyle{tcl}{
  style=base,
  language=tcl,
  morekeywords={read_vhdl, read_xdc, synth_design, opt_design, place_design, route_design,
    write_bitstream, report_timing_summary, report_utilization, report_timing,
    open_hw_manager, connect_hw_server, open_hw_target, program_hw_devices},
}

% Ambienti per il codice scritto direttamente nel .tex
\lstnewenvironment{vhdl}[1][]{\lstset{style=vhdl,#1}}{}
\lstnewenvironment{xdc}[1][]{\lstset{style=xdc,#1}}{}
\lstnewenvironment{tcl}[1][]{\lstset{style=tcl,#1}}{}

% Cartella in cui si trovano i sorgenti da includere (relativa a main.tex)
\newcommand{\codedir}{codice}

% Inclusione di sorgenti da file: il codice nell'elaborato resta allineato a quello simulato.
%   \vhdlfile[firstline=10, lastline=30]{es1/mux16.vhd}{Mux 16:1 strutturale}{lst:mux16}
\newcommand{\vhdlfile}[4][]{\lstinputlisting[style=vhdl, caption={#3}, label={#4}, #1]{\codedir/#2}}
\newcommand{\xdcfile}[4][]{\lstinputlisting[style=xdc, caption={#3}, label={#4}, #1]{\codedir/#2}}

% Nomi di segnali, porte ed entity nel testo
\newcommand{\sig}[1]{\texttt{#1}}
\newcommand{\ent}[1]{\texttt{\textbf{#1}}}

%----------------------------------------------------------------------------------------
%	RIQUADRI DELL'ELABORATO
%----------------------------------------------------------------------------------------

% Stile comune: barra verticale colorata a sinistra, sfondo tenue, titolo in grassetto.
\newcommand{\definisciriquadro}[3]{% #1 = nome stile, #2 = colore, #3 = sfondo (% di colore)
	\mdfdefinestyle{#1}{
		topline=false, bottomline=false, rightline=false,
		linewidth=3pt, linecolor=#2,
		backgroundcolor=#2!#3!white,
		innerleftmargin=10pt, innerrightmargin=10pt,
		innertopmargin=7pt, innerbottommargin=7pt,
		leftmargin=0pt, rightmargin=0pt,
		skipabove=\medskipamount, skipbelow=\medskipamount,
		startcode={\setlength{\parindent}{0pt}\setlength{\parskip}{0.4em}},
	}%
}

\definisciriquadro{stile-traccia}{traccol}{6}
\definisciriquadro{stile-scelta}{accento}{6}
\definisciriquadro{stile-osservazione}{osservcol}{6}
\definisciriquadro{stile-attenzione}{attcol}{6}
\definisciriquadro{stile-risultato}{riscol}{5}

% Usage: \begin{traccia} testo ufficiale dell'esercizio \end{traccia}
\newenvironment{traccia}
	{\begin{mdframed}[style=stile-traccia]\noindent{\sffamily\bfseries\color{traccol}Traccia}\par\smallskip\itshape}
	{\end{mdframed}}

% Scelte progettuali numerate per capitolo, citabili con \ref
% Usage: \begin{scelta}[Titolo breve]\label{sc:...} ... \end{scelta}
\newcounter{scelta}[chapter]
\renewcommand{\thescelta}{\thechapter.\arabic{scelta}}
\newenvironment{scelta}[1][]
	{\refstepcounter{scelta}\begin{mdframed}[style=stile-scelta]%
	 \noindent{\sffamily\bfseries\color{accento}Scelta progettuale \thescelta\if\relax\detokenize{#1}\relax\else\ --- #1\fi}\par\smallskip}
	{\end{mdframed}}

% Variante non numerata: \begin{scelta*}[Titolo] ... \end{scelta*}
\newenvironment{scelta*}[1][]
	{\begin{mdframed}[style=stile-scelta]%
	 \noindent{\sffamily\bfseries\color{accento}Scelta progettuale\if\relax\detokenize{#1}\relax\else\ --- #1\fi}\par\smallskip}
	{\end{mdframed}}

% Usage: \begin{osservazione}[titolo opzionale] ... \end{osservazione}
\newenvironment{osservazione}[1][Osservazione]
	{\begin{mdframed}[style=stile-osservazione]\noindent{\sffamily\bfseries\color{osservcol}#1}\par\smallskip}
	{\end{mdframed}}

% Usage: \begin{attenzione}[titolo opzionale] ... \end{attenzione}
\newenvironment{attenzione}[1][Attenzione]
	{\begin{mdframed}[style=stile-attenzione]\noindent{\sffamily\bfseries\color{attcol}#1}\par\smallskip}
	{\end{mdframed}}

% Esito di una simulazione / prova su board
% Usage: \begin{risultato}[titolo opzionale] ... \end{risultato}
\newenvironment{risultato}[1][Risultato]
	{\begin{mdframed}[style=stile-risultato]\noindent{\sffamily\bfseries\color{riscol}#1}\par\smallskip}
	{\end{mdframed}}

%----------------------------------------------------------------------------------------
%	TABELLA DELLE PORTE (interfaccia di una entity)
%----------------------------------------------------------------------------------------

% Usage:
% \begin{porte}[Interfaccia di \ent{mux16}]
%   \porta{d}{in}{16}{dati in ingresso}
% \end{porte}
\newcolumntype{Y}{>{\raggedright\arraybackslash}X}
\newenvironment{porte}[1][]
	{\begin{table}[H]\centering\small
	 \if\relax\detokenize{#1}\relax\else\caption{#1}\fi
	 \begin{tabular}{@{}l c c >{\raggedright\arraybackslash}p{0.58\linewidth}@{}}\toprule
	 \textbf{Porta} & \textbf{Direzione} & \textbf{Bit} & \textbf{Descrizione} \\ \midrule}
	{\bottomrule\end{tabular}\end{table}}
\newcommand{\porta}[4]{\texttt{#1} & \texttt{#2} & #3 & #4 \\}

%----------------------------------------------------------------------------------------
%	SCREENSHOT E IMMAGINI
%----------------------------------------------------------------------------------------

% Inserisce un'immagine se il file esiste, altrimenti un segnaposto con il nome del file
% atteso: il documento compila anche prima di aver fatto gli screenshot.
% Usage: \screenshot[0.9\linewidth]{sim_mux16.png}{Simulazione del mux 16:1}{fig:sim-mux16}
\newcommand{\screenshot}[4][\linewidth]{%
	\begin{figure}[H]\centering
	\IfFileExists{figure/screenshot/#2}{%
		\fbox{\includegraphics[width=#1]{figure/screenshot/#2}}%
	}{%
		\begin{tikzpicture}
			\node[draw=black!40, dashed, fill=black!3, minimum width=#1, minimum height=4.5cm,
			      align=center, font=\sffamily\small\color{black!55}]
			     {Screenshot da inserire\\[2pt]\texttt{figure/screenshot/\detokenize{#2}}};
		\end{tikzpicture}%
	}
	\caption{#3}\label{#4}
	\end{figure}}

%----------------------------------------------------------------------------------------
%	STILI TIKZ PER SCHEMI A BLOCCHI, FSM E DIAGRAMMI TEMPORALI
%----------------------------------------------------------------------------------------

\tikzset{
	>={Stealth[length=2mm]},
	% blocco generico (componente)
	blocco/.style={draw, thick, rectangle, fill=accento!5, minimum width=2cm, minimum height=1cm,
	               align=center, font=\sffamily\small},
	% multiplexer (trapezio)
	mux/.style={draw, thick, trapezium, trapezium angle=70, shape border rotate=270,
	            fill=accento!5, minimum width=1.6cm, minimum height=0.8cm, font=\sffamily\scriptsize,
	            inner sep=1pt},
	% registro
	reg/.style={draw, thick, rectangle, fill=osservcol!8, minimum width=1.4cm, minimum height=0.8cm,
	            font=\sffamily\small},
	% unità di controllo
	ctrl/.style={draw, thick, rounded corners=4pt, fill=attcol!6, minimum width=2.4cm, minimum height=1cm,
	             align=center, font=\sffamily\small},
	% stato di una FSM
	fsmstate/.style={state, thick, fill=accento!5, minimum size=1.2cm, font=\sffamily\small, align=center},
	% bus (più bit)
	bus/.style={->, very thick},
	% segnale di controllo
	seg/.style={->, thin},
	% barretta con larghezza del bus:  \draw[bus] (a) -- node[larghezza=8]{/} (b);
	larghezza/.style={inner sep=0pt, font=\small,
	                  label={[font=\scriptsize, inner sep=1pt]above:#1}},
	% diagrammi temporali
	onda/.style={thick, accento},
	grigliatempo/.style={black!15, thin},
}

% Numero cerchiato per le fasi di un protocollo, utilizzabile anche nelle didascalie: \fase{1}
\DeclareRobustCommand{\fase}[1]{\tikz[baseline=(c.base)]\node[circle, draw, inner sep=0.6pt, font=\scriptsize\sffamily](c){#1};}

%----------------------------------------------------------------------------------------
%	COMMAND LINE ENVIRONMENT
%----------------------------------------------------------------------------------------

% Usage:
% \begin{commandline}
%	\begin{verbatim}
%		$ ls
%
%		Applications	Desktop	...
%	\end{verbatim}
% \end{commandline}

\mdfdefinestyle{commandline}{
	leftmargin=10pt,
	rightmargin=10pt,
	innerleftmargin=15pt,
	middlelinecolor=black!50!white,
	middlelinewidth=2pt,
	frametitlerule=false,
	backgroundcolor=black!5!white,
	frametitle={Command Line},
	frametitlefont={\normalfont\sffamily\color{white}\hspace{-1em}},
	frametitlebackgroundcolor=black!50!white,
	nobreak,
}

% Define a custom environment for command-line snapshots
\newenvironment{commandline}{
	\medskip
	\begin{mdframed}[style=commandline]
}{
	\end{mdframed}
	\medskip
}

%----------------------------------------------------------------------------------------
%	FILE CONTENTS ENVIRONMENT
%----------------------------------------------------------------------------------------

% Usage:
% \begin{file}[optional filename, defaults to "File"]
%	File contents, for example, with a listings environment
% \end{file}

\mdfdefinestyle{file}{
    innertopmargin=1.6\baselineskip,
    innerbottommargin=0.8\baselineskip,
    topline=false, bottomline=false,
    leftline=false, rightline=false,
    leftmargin=2cm,
    rightmargin=2cm,
    singleextra={
        \draw[fill=black!10!white](P)++(0,-1.2em)rectangle(P-|O);
        \node[anchor=north west] at(P-|O){\ttfamily\mdfilename};
        \def\l{3em}
        \draw(O-|P)++(-\l,0)--++(\l,\l)--(P)--(P-|O)--(O)--cycle;
        \draw(O-|P)++(-\l,0)--++(0,\l)--++(\l,0);
    },
}

% Comando per tab a 4 spazi all’interno dell’ambiente file.
% Il tab attivo va definito in un gruppo con il catcode già cambiato, altrimenti
% \def^^I verrebbe letto come \def seguito da uno spazio.
\begingroup
\catcode`\^^I=\active
\gdef\attivatab{\catcode`\^^I=\active\def^^I{\hspace*{4ex}}}
\endgroup

\newenvironment{file}[1][File]{%
    \medskip
    \newcommand{\mdfilename}{#1}
    \attivatab
    \begin{mdframed}[style=file]
}{%
    \end{mdframed}
    \medskip
}

%----------------------------------------------------------------------------------------
%	NUMBERED QUESTIONS ENVIRONMENT
%----------------------------------------------------------------------------------------

% Usage:
% \begin{question}[optional title]
%	Question contents
% \end{question}

\mdfdefinestyle{question}{
	innertopmargin=1.2\baselineskip,
	innerbottommargin=0.8\baselineskip,
	roundcorner=5pt,
	nobreak,
	singleextra={%
		\draw(P-|O)node[xshift=1em,anchor=west,fill=white,draw,rounded corners=5pt]{%
		Domanda \theQuestion\questionTitle};
	},
}

\newcounter{Question} % Stores the current question number that gets iterated with each new question

% Define a custom environment for numbered questions
\newenvironment{question}[1][\unskip]{
	\bigskip
	\stepcounter{Question}
	\newcommand{\questionTitle}{~#1}
	\begin{mdframed}[style=question]
}{
	\end{mdframed}
	\medskip
}

%----------------------------------------------------------------------------------------
%	WARNING TEXT ENVIRONMENT
%----------------------------------------------------------------------------------------

% Usage:
% \begin{warn}[optional title, defaults to "Warning:"]
%	Contents
% \end{warn}

\mdfdefinestyle{warning}{
	topline=false, bottomline=false,
	leftline=false, rightline=false,
	nobreak,
	singleextra={%
		\draw(P-|O)++(-0.5em,0)node(tmp1){};
		\draw(P-|O)++(0.5em,0)node(tmp2){};
		\fill[black,rotate around={45:(P-|O)}](tmp1)rectangle(tmp2);
		\node at(P-|O){\color{white}\scriptsize\bf !};
		\draw[very thick](P-|O)++(0,-1em)--(O);%--(O-|P);
	}
}

% Define a custom environment for warning text
\newenvironment{warn}[1][Attenzione:]{
	\medskip
	\begin{mdframed}[style=warning]
		\noindent{\textbf{#1}}
}{
	\end{mdframed}
}

%----------------------------------------------------------------------------------------
%	INFORMATION ENVIRONMENT
%----------------------------------------------------------------------------------------

% Usage:
% \begin{info}[optional title, defaults to "Info:"]
% 	contents
% 	\end{info}

\mdfdefinestyle{info}{%
	topline=false, bottomline=false,
	leftline=false, rightline=false,
	nobreak,
	singleextra={%
		\fill[black](P-|O)circle[radius=0.4em];
		\node at(P-|O){\color{white}\scriptsize\bf i};
		\draw[very thick](P-|O)++(0,-0.8em)--(O);%--(O-|P);
	}
}

% Define a custom environment for information
\newenvironment{info}[1][Info:]{
	\medskip
	\begin{mdframed}[style=info]
		\noindent{\textbf{#1}}
}{
	\end{mdframed}
}
